\subsubsection*{Solution}
\paragraph*{Step-by-Step Derivation}

Let the plate separation be $d$, and nondimensionalize

\[
z=\frac{z^\ast}{d},\qquad t=\frac{\kappa t^\ast}{d^2},\qquad T=\frac{T^\ast-T^\ast_{\rm top}}{q d/K},
\]

where $q$ is the imposed upward heat flux and $K$ is the thermal conductivity. The corresponding flux-based Rayleigh number is

\[
Ra=\frac{g\alpha qd^4}{K\nu\kappa} =\frac{g\alpha \Delta T_{\rm cond}d^3}{\nu\kappa}, \qquad \Delta T_{\rm cond}=\frac{qd}{K}.
\]

\textcolor{blue}{where $\alpha$ is thermal expansion coefficient, $\nu$ is viscosity, $\kappa$ is thermal diffusivity.} Thus, at the conduction state, the flux-based and temperature-drop-based Rayleigh numbers coincide. Fixed flux is a Neumann thermal condition, while fixed temperature is Dirichlet; this normalization and distinction are standard for mixed thermal RBC \href{\detokenize{https://www.sfu.ca/~rwwitten/papers/WittenbergGao_epjb2010.pdf}}{Wittenberg and Gao}. The bottom-flux/top-temperature configuration is commonly called an "FT" configuration \href{\detokenize{https://arxiv.org/abs/2003.00026}}{Anders et al.}.

\subparagraph*{1. Conduction state and perturbation equations}

The nondimensional Boussinesq equations are

\[
Pr^{-1}\left(\frac{\partial\mathbf u}{\partial t} +\mathbf u\cdot\nabla\mathbf u\right) =-\nabla p+\nabla^2\mathbf u+Ra\,T\,\hat{\mathbf z},
\]

\[
\frac{\partial T}{\partial t}+\mathbf u\cdot\nabla T=\nabla^2T, \qquad \nabla\cdot\mathbf u=0.
\]

The conductive base state is

\[
\mathbf u_0=0,\qquad T_0(z)=1-z,
\]

because

\[
-DT_0(0)=1,\qquad T_0(1)=0,
\]

where $D=d/dz$.

Write

\[
T=T_0+\theta,\qquad \mathbf u=\mathbf u',
\]

and retain terms linear in the perturbations. For a horizontal Fourier mode,

\[
(w,\theta) = \bigl(W(z),\Theta(z)\bigr) e^{\sigma t+i\mathbf a\cdot\mathbf x}, \qquad a=|\mathbf a|,
\]

define

\[
L=D^2-a^2.
\]

Eliminating pressure gives

\[
Pr^{-1}\sigma LW=L^2W-Ra\,a^2\Theta,
\]

while the temperature equation becomes

\[
(\sigma-L)\Theta=W.
\]

\subparagraph*{2. The marginal mode is stationary}

To rule out an oscillatory first instability, let

\[
U=\int|\mathbf u|^2\,dV,\qquad \Theta_2=\int|\theta|^2\,dV,
\]

and define the positive dissipations

\[
A=\int|\nabla\mathbf u|^2\,dV,\qquad B=\int|\nabla\theta|^2\,dV.
\]

Taking inner products of the linear momentum and heat equations gives

\[
Pr^{-1}\sigma U=-A+Ra\,I^\ast,
\]

\[
\sigma\Theta_2=-B+I, \qquad I=\int w\theta^\ast\,dV.
\]

Thus

\[
I=B+\sigma\Theta_2.
\]

If $\sigma=s+i\omega$, the imaginary part of the momentum identity is

\[
\frac{\omega}{Pr}U=-Ra\,\omega\Theta_2,
\]

or

\[
\omega\left(\frac{U}{Pr}+Ra\,\Theta_2\right)=0.
\]

Since the quantity in parentheses is positive, $\omega=0$. Therefore marginal stability occurs at

\[
\sigma=0.
\]

Consequently,

\[
L^2W=Ra\,a^2\Theta,\qquad L\Theta=-W.
\]

Eliminating $\Theta$ yields the sixth-order eigenvalue equation

\[
{L^3W=-Ra\,a^2W}.
\]

Equivalently,

\[
W^{(6)}-3a^2W^{(4)}+3a^4W'' +\left(Ra\,a^2-a^6\right)W=0.
\]

\subparagraph*{3. Boundary conditions}

At the no-slip, constant-flux bottom boundary $z=0$,

\[
W=0,\qquad DW=0,\qquad D\Theta=0.
\]

At the impermeable free-slip, fixed-temperature top boundary $z=1$,

\[
W=0,\qquad D^2W=0,\qquad \Theta=0.
\]

Using

\[
\Theta=\frac{L^2W}{Ra\,a^2},
\]

the thermal conditions become

\[
D(L^2W)=0\quad(z=0), \qquad L^2W=0\quad(z=1).
\]

Hence the six boundary conditions on $W$ are

\[
W(0)=W'(0)=0, \qquad W^{(5)}(0)-2a^2W^{(3)}(0)=0,
\]

\[
W(1)=W''(1)=W^{(4)}(1)=0.
\]

\subparagraph*{4. Numerical neutral curve}

Introducing

\[
\mathbf y=(W,W',W'',W^{(3)},W^{(4)},W^{(5)})^T,
\]

the eigenproblem can be written as $\mathbf y'=A\mathbf y$, with

\[
A=
\begin{pmatrix}
0&1&0&0&0&0\\
0&0&1&0&0&0\\
0&0&0&1&0&0\\
0&0&0&0&1&0\\
0&0&0&0&0&1\\
a^6-Ra\,a^2&0&-3a^4&0&3a^2&0
\end{pmatrix}.
\]

A basis satisfying the three bottom conditions is

\[
Y_0=
\begin{pmatrix}
0&0&0\\
0&0&0\\
1&0&0\\
0&1&0\\
0&0&1\\
0&2a^2&0
\end{pmatrix}.
\]

The three top conditions require

\[
F(Ra,a)
=
\det\!\left[
\begin{pmatrix}
1&0&0&0&0&0\\
0&0&1&0&0&0\\
0&0&0&0&1&0
\end{pmatrix}
e^A Y_0
\right]=0.
\]

For each $a$, the smallest positive root of $F=0$ gives the neutral curve $Ra_0(a)$. Minimizing that curve gives

\[
a_c=2.2147266\ldots,
\]

\[
Ra_c=816.7444317\ldots.
\]

The corresponding horizontal wavelength is

\[
\lambda_c=\frac{2\pi}{a_c}=2.8370028\ldots\,d.
\]

As an external check, the same boundary combination-fixed temperature and stress-free at one boundary, thermally insulating perturbation and no-slip at the other-is tabulated in the literature as $Ra_c=816.748$ and $a_c=2.21$ \href{\detokenize{https://www.zimmermann.physik.uni-bayreuth.de/pool/publikationen/2009/Homologous_onset_of_double_layer_convection/Homologous_onset_of_double_layer_convection.pdf}}{Busse and Petry, Eq. (32)}.
